\par\indent\SourceLeaf{229}%
Methodus \textit{secunda} sistit coordinatas in forma functionum duarum variabilium \(p,q\). Supponamus per differentiationem harum functionum prodire
\[
\begin{aligned}
\mathrm{d}x&=a\mathrm{d}p+a'\mathrm{d}q\\
\mathrm{d}y&=b\mathrm{d}p+b'\mathrm{d}q\\
\mathrm{d}z&=c\mathrm{d}p+c'\mathrm{d}q
\end{aligned}
\]
\noindent quibus valoribus in formula supra data substitutis, obtinemus
\[
(aX+bY+cZ)\mathrm{d}p+(a'X+b'Y+c'Z)\mathrm{d}q=0
\]
\noindent Quum haec aequatio locum habere debeat independenter a valoribus differentialium \(\mathrm{d}p,\mathrm{d}q\), manifesto esse debebit
\[
aX+bY+cZ=0,\qquad a'X+b'Y+c'Z=0
\]
\noindent unde colligimus, \(X,Y,Z\) proportionales esse debere quantitatibus
\[
bc'-cb',\quad ca'-ac',\quad ab'-ba'
\]
\par\indent Statuendo itaque brevitatis caussa
\[
\BareRoot{((bc'-cb')^2+(ca'-ac')^2+(ab'-ba')^2)}=\Delta
\]
\noindent erit vel
\[
X=\frac{bc'-cb'}{\Delta},\quad Y=\frac{ca'-ac'}{\Delta},\quad Z=\frac{ab'-ba'}{\Delta}
\]
\noindent vel
\[
X=\frac{cb'-bc'}{\Delta},\quad Y=\frac{ac'-ca'}{\Delta},\quad Z=\frac{ba'-ab'}{\Delta}
\]
\par\indent His duabus methodis generalibus accedit \textit{tertia}, ubi una coordinatarum, e. g. \(z\) exhibetur in forma functionis reliquarum \(x,y\): haec methodus manifesto nihil aliud est, nisi casus specialis vel methodi primae, vel secundae. Quodsi hic statuitur
\[
\mathrm{d}z=t\mathrm{d}x+u\mathrm{d}y
\]
\noindent erit vel
\[
X=\frac{-t}{\BareRoot{(1+tt+uu)}},\quad Y=\frac{-u}{\BareRoot{(1+tt+uu)}},\quad Z=\frac{1}{\BareRoot{(1+tt+uu)}}
\]
\noindent vel
\[
X=\frac{t}{\BareRoot{(1+tt+uu)}},\quad Y=\frac{u}{\BareRoot{(1+tt+uu)}},\quad Z=\frac{-1}{\BareRoot{(1+tt+uu)}}
\]
